% Part: incompleteness
% Chapter: introduction
% Section: historical-background

\documentclass[../../../include/open-logic-section]{subfiles}

\begin{document}

\olfileid{inc}{int}{bgr}

\olsection{Latar Sejarah}

Ing bagean iki, kita bakal ngrembug kanthi ringkes
perkembangan sejarah kang mbantu manggonake
teorema-teorema ngenani ora jangkep ing konteks.
Mligine, kita bakal menehi gegambaran garis gedhé
banget ngenani sejarah logika matematika, banjur
ngandharake sawatara bab ngenani sejarah
landhesan matematika.

\begin{digress}
  Frasa ``logika matematika'' nduwé teges loro.
  Tembung ``matematika'' bisa ditegesi njlentrehake
  bab kang disinaoni, kaya ing ``logikane matematika'',
  yaiku prinsip-prinsip nalar matematika; utawa
  njlentrehake carane nyinaoni, kaya ing
  ``matematikane logika'', yaiku pangkajian matematis
  marang prinsip-prinsip nalar. Andharan ing ngisor
  iki nyakup logika matematika ing rong teges mau,
  asring bebarengan.
\end{digress}

Pangkajian logika wiwit, satemene, saka Aristoteles,
kang urip kira-kira taun 384--322 \textsc{sm}.
Karyane \emph{Categories}, \emph{Prior
  analytics}, lan \emph{Posterior analytics} ngemot
pangkajian sistematis marang prinsip-prinsip nalar
ilmiah, kalebu pangkajian silogisme kang jero lan
sistematis.

Logikane Aristoteles dadi kang utama ing filsafat
skolastik nganti Abad Pertengahan; malah ing abad
kaping wolulas, Kant isih nganggep logika Aristoteles
wis sampurna lan ora perlu direvisi. Nanging teori
silogisme winates banget, saengga mung bisa
nggambarake pérangan kang paling lumahing saka
nalar matematika. Sakabad sadurunge, Leibniz,
kanca sezamane Newton, mbayangake ``kalkulus''
kang jangkep kanggo nalar logis, lan wis miwiti
langkah-langkah prasaja kanggo ngrancang kalkulus
kaya mangkono, kang pokoke njlentrehake salah
sijining wujud logika proposisional.

Abad kaping sangalas dadi titik owah wigati kanggo
logika. Ing 1854 George Boole nulis \emph{The Laws of
  Thought}, kang ngemot pangkajian aljabar logika
proposisional kanthi jero, ora adoh saka panyajian
modern. Ing 1879 Gottlob Frege nerbitake
\emph{Begriffsschrift} (tulisan konsep), kang
ngembangake logika proposisional nganggo kuantor
lan relasi, mula uga nyakup logika tataran kapisan.
Sejatine, sistem logikane Frege uga nyakup logika
tataran luwih dhuwur lan liya-liyane. Ing \emph{Basic Laws
  of Arithmetic}, Frege kepengin nuduhake yen
kabeh aritmetika bisa didudut ing Begriffsschrift
saka asumsi logis murni. Sayange, asumsi-asumsi
iki jebul ora konsisten, kaya kang dituduhake
Russell ing 1902. Nanging yen aksioma kang ora
konsisten mau disisihake, Frege meh bisa diarani
nemokake logika modern kanthi upayane dhewe,
sawijining prestasi kang nggumunake. Logika
kuantifikasi uga dikembangake sacara mandiri
dening para pamikir kang cenderung aljabar sawisé
Boole, kalebu Peirce lan Schr\"oder.

Saiki ayo ngalih menyang perkembangan landhesan
matematika. Mesthi wae, amarga logika nduwé peran
wigati ing matematika, perkembangan iki akeh
sesambungane karo kang mau diandharake. Tuladhane,
Frege ngembangake logikane kanthi ancas cetha
kanggo nuduhake yen kabeh matematika bisa
adhedhasar kerangka logikane wae. Mligine, dheweke
kepengin nuduhake yen matematika dumadi saka
kabeneran a priori kang \emph{analitis}, dudu
kabeneran a priori kang \emph{sintetis} kaya
panemune Kant.

Akeh wong nganggep matematika sajatine wiwit
ing bangsa Yunani. \emph{Elements} anggitane
Euklides, ditulis kira-kira taun 300 SM, wis dadi
conto matematika Yunani kang mateng, kanthi
tekanan marang katelitian lan kecermatan.
Definisi lan bukti ing \emph{Elements} anggitane
Euklides isih lestari meh tanpa owah ing buku
geometri sekolah menengah saiki (nalika geometri
isih diwulangake ing sekolah menengah). Pola
nalar matematika iki dianggep tuladha kanggo
argumentasi kang tliti, ora mung ing matematika
nanging uga ing cabang-cabang filsafat. (Spinoza
malah nyajekake argumentasi moral lan agama
nganggo gaya Euklides, kang katon aneh!)

Kalkulus ditemokake dening Newton lan Leibniz ing
abad kaping pitulas. (Padu ngenani sapa kang luwih
dhisik lumaku nganti pirang-pirang abad, nanging
akeh sarjana saiki nganggep loro perkembangan iku
umume mandiri.) Kalkulus nglibatake nalar ngenani,
umpamane, jumlah tanpa wates saka besaran
infinitesimal; ciri-ciri iki nyurung
kritik saka Uskup Berkeley, kang negesake yen
kapercayan marang Gusti Allah ora kurang rasional
tinimbang matematika ing jamane. Cara-cara
kalkulus digunakake kanthi amba ing abad kaping
wolulas, umpamane dening Leonhard Euler, kang
nggunakake petungan jumlah tanpa wates lan
ngasilake asil kang nggumunake.

Ing abad kaping sangalas, para matematikawan
nyoba mangsuli kritik Berkeley kanthi menehi
landhesan kang luwih kukuh marang kalkulus.
Upayane Cauchy, Weierstrass, Bolzano, lan liya-liyane
ngasilake definisi saiki kanggo limit, kontinuitas,
diferensiasi, lan integrasi nganggo ``epsilon lan
delta'', yaiku tanpa nyebut besaran infinitesimal.
Ing pungkasan abad iku, para matematikawan nyoba
luwih adoh maneh: njlentrehake kabeh pérangan
kalkulus, kalebu wilangan real, adhedhasar wilangan
natural. (Kronecker: ``Gusti Allah nitahake wilangan
wutuh; liyane kabeh gaweyane manungsa.'') Ing 1872,
Dedekind nulis ``Kontinuitas lan wilangan irasional'',
kang nuduhake cara ``mbangun'' wilangan real minangka
himpunan wilangan rasional (kang, kaya wis kok ngerti,
bisa dipandang minangka pasangan wilangan natural).
Ing 1888 dheweke nulis ``Was sind und was sollen die
Zahlen'' (kira-kira, ``Apa iku wilangan natural, lan
kepriye kudune?'') kanggo njlentrehake wilangan
natural kanthi istilah kang murni ``logis''. Ing 1887
Kronecker nulis ``\"Uber den Zahlbegriff''
(``Ngenani konsep wilangan''), kang ngrembug cara
nggambarake kabeh objek matematika nganggo
wilangan wutuh; ing 1889 Giuseppe Peano menehi
aksioma formal lan simbolis kanggo wilangan natural.

Pungkasan abad kaping sangalas uga nggawa
wani anyar nalika ngadhepi tanpa wates.
Sadurunge iku, objek lan struktur tanpa wates
(kayata himpunan wilangan natural) ditangani kanthi
ati-ati; ``tanpa wates akehe'' dimangerteni
minangka ``sakarepmu akehe'', lan ``nyedhaki
limit'' dimangerteni minangka ``cedhak sakarepmu''.
Nanging Georg Cantor nuduhake yen tanpa wates
bisa ditampa miturut teges harfiahe. Pakaryane
Cantor, Dedekind, lan liya-liyane mbantu
ngenalake cara mangerteni matematika liwat teori
himpunan kang saiki ditampa kanthi amba.

Iki nggawa kita menyang perkembangan logika lan
landhesan matematika ing abad kaping rong puluh.
Ing 1902 Russell nemokake paradoks ing sistem
logikane Frege. Ing 1904 Zermelo mbuktekake
prinsip urutan rapi Cantor nganggo kang diarani
``aksioma pilihan''; sah-orane aksioma iki
nuwuhake debat kang akeh. Antarane 1910 lan 1913,
telung jilid \emph{Principia Mathematica} anggitane
Russell lan Whitehead terbit, nerusake program
Frege kanggo nglandhesake matematika ing logika.
Sayange, Russell lan Whitehead kudu nggunakake
rong prinsip kang angel dibenerake minangka
prinsip logis murni: aksioma tanpa wates lan
aksioma ``reduksibilitas''. Ing dasawarsa 1900-an,
Poincar\'e ngritik panggunaan ``definisi
impredikatif'' ing matematika, lan ing dasawarsa
1910-an Brouwer wiwit ngusulake supaya kabeh
matematika diwangun maneh ing landhesan
``intuisionistis'' kang ora nggunakake hukum
ora anane pilihan katelu ($!A \lor \lnot !A$).

Pancen jaman kang aneh!{} Program kanggo nyuda
kabeh matematika dadi logika saiki diarani
``logisisme'', lan umume dianggep gagal amarga
kangelan-kangelan kang kasebut ing ndhuwur.
Program kanggo ngembangake matematika adhedhasar
konstruksi mental intuisionistis diarani
``intuisionisme'', lan dianggep menehi watesan
kang banget ketat kanggo matematika saben dina.
Ing watara wiwitan abad iku, David Hilbert,
salah sawijining matematikawan kang paling
duwe pengaruh ing sejarah, dadi panyengkuyung
kuwat cara-cara abstrak anyar saka Cantor lan
Dedekind: ``ora ana wong kang bisa ngusir kita
saka swarga kang digawe Cantor kanggo kita''.
Ing wektu kang padha, dheweke uga nggatekake
kritik landhesan marang cara-cara anyar iki
(kang anehé, saiki diarani ``klasik'').
Dheweke ngusulake cara kanggo nggayuh rong
ancas bebarengan:
\begin{enumerate}
\item Gambarake cara klasik nganggo aksioma lan
  aturan formal; gambarake pitakon matematika
  minangka !!{formula}s ing sistem aksiomatik.
\item Gunakna cara kang aman lan ``finitèr''
  kanggo mbuktekake yen sistem deduktif formal
  iki konsisten.
\end{enumerate}

Pakaryane Hilbert maju adoh kanggo nggayuh ancas
kapisan. Ing 1899, dheweke wis nindakake iki
kanggo geometri ing bukune kang misuwur,
\emph{Foundations of geometry}. Ing taun-taun
sabanjure, dheweke bebarengan karo para murid
lan kanca kerjane ngupaya nindakake bab kang padha
ing cabang matematika liyane. Hilbert dhewe
menehi sistem aksioma kanggo aritmetika lan analisis.
Zermelo menehi aksiomatisasi teori himpunan, kang
banjur dikembangake dening Fraenkel, Skolem,
von Neumann, lan liya-liyane. Ing tengah
dasawarsa 1920-an, ana rong pendekatan kang
ngaku minangka aksiomatisasi kanggo ``kabeh''
matematika: \emph{Principia mathematica}
anggitane Russell lan Whitehead, lan kang banjur
dikenal minangka teori himpunan Zermelo--Fraenkel.

Ing 1921, Hilbert miwiti proyek riset kanggo
mbuktekake yen sistem-sistem iki konsisten.
Dheweke dibantu para muride, mligine Bernays,
Ackermann, lan mengko Gentzen. Gagasan pokok
kanggo nggayuh ancas iki yaiku ngowahi pitakon
apa inkonsistensi ing matematika bisa didudut
liwat !!{derivation} dadi masalah kombinatorial
ngenani urutan simbol kang mungkin. Yaiku,
urutan ukara kang mungkin lan nyukupi kritéria
minangka !!{derivation} kang bener tumrap,
umpamane, $!A \land \lnot !A$ saka aksioma
sistem aksiomatik kanggo aritmetika, analisis,
utawa teori himpunan. Bukti yen urutan simbol
kaya mangkono ora mungkin ana---amarga bukti iku
dhewe bukti matematis---dikarepake bisa diformalkan ing
sistem aksiomatik kasebut. Kanthi tembung liya,
ana sawijining ukara $\OCon$ kang nyatakake,
umpamane, aritmetika konsisten. Kajaba iku,
ukara iki mesthine bisa dibuktekake ing sistem
kang dirembug, mligine yen buktine mung mbutuhake
cara kang winates banget, utawa ``finitèr''.

Ancas kapindho, yaiku supaya sistem aksioma kang
dikembangake bisa mutusake saben pitakon matematika,
bisa ditegesi kanthi cetha nganggo rong cara.
Cara kapisan mangkene: Kanggo saben ukara~$!A$
ing basa sawijining sistem aksioma matematika,
$!A$ utawa $\lnot !A$ bisa dibuktekake saka
aksioma-aksiomane. Yen iki bener, ora bakal ana
ukara kang ora bisa dibuktekake lan uga ora bisa
dibantah adhedhasar aksioma, ora ana pitakon
kang ora bisa diputusake dening aksioma. Sistem
aksioma kang nduwé sifat iki diarani \emph{jangkep}.
Mesthi wae, kanggo ukara tartamtu bisa uga angel
nemtokake endi saka rong alternatif iku kang bener.
Nanging miturut prinsip, kudune ana cara kanggo
nindakake iku. Malah kanggo sistem aksioma lan
!!{derivation} kang dirembug Hilbert, kajangkepan
bakal nyebabake anane cara kaya mangkono---sanajan
Hilbert ora ngerti bab iki. Cara kapindho kanggo
negesi pitakon iku yaiku syarat kang luwih kuwat:
kudu ana cara mekanis utawa komputasional kang
bisa nemtokake, kanggo ukara~$!A$ tartamtu,
apa ukara iku bisa didudut saka aksioma utawa ora.

Ing 1931, G\"odel mbuktekake rong ``teorema ngenani
ora jangkep'', kang nuduhake yen program iki ora
bisa kasil.
Cathetan koreksi sumber OLPL-166 lan SOL6-F068: ringkesan ing ngisor iki dibenerake supaya nyebut syarat teori kang konsisten, efektif diaksiomatisasi lan cukup kuwat. Ukara konsistensi ora kanthi otomatis dijamin ora bisa dibantah dening teorema kapindho.
Sistem aksioma kanggo matematika kang konsisten, efektif
diaksiomatisasi lan cukup kuwat ora jangkep. Miturut
syarat teorema kapindho kang dijlentrehake mengko,
ukara kang nyatakake konsistensine dhewe ora bisa
dibuktekake saka aksioma sistem iku.

Iki dadi pukulan abot marang program asli Hilbert.
Nanging, kaya kang asring kelakon ing matematika,
iki uga mbukak dalan riset anyar kang narik kawigaten.
Yen ora ana siji sistem formal matematika kang
nyakup kabeh, kita bisa ngembangake sistem kang
luwih winates cakupane lan nliti apa kang bisa
dibuktekake ing kono. Kita uga bisa ngembangake
cara bukti kang ora ketat watesane kanggo
netepake konsistensi sistem-sistem mau, lan
nggoleki cara kanggo ngukur sepira angele
mbuktekake konsistensine. Amarga G\"odel nuduhake
yen sistem formal kang konsisten, efektif diaksiomatisasi
lan cukup kuwat nduwé pitakon
kang ora bisa diputusake, kita bisa nggoleki
pitakon kang ``narik kawigaten'' lan ora bisa
diputusake dening sawijining sistem formal,
banjur nemtokake sepira kuate sistem formal
kang dibutuhake kanggo mutusake pitakon iku.
Nganti saiki, para ahli logika terus nliti
pitakon-pitakon iki ing disiplin matematika
anyar, yaiku teori bukti.

\end{document}
